\chapter{Conclusion}\label{ch:conclusion}

\section{Summary of the Research}\label{sec:research_summary}

This work investigated how deep reinforcement learning agents internally represent and pursue sequential goals in the procedurally generated Procgen Heist environment. Using mechanistic interpretability techniques such as linear probes, activation interventions, parallel rollout analysis, and SAEs, we investigated the core mechanism behind how our network learned to encode goal structures in this multi-objective task and whether interpretability tools developed for language models translate to RL agents.

\section{Addressing Research Questions and Hypotheses}
\label{sec:addressing_rq}

Before interpreting our findings in detail, we first highlight the original research questions and hypotheses from Chapter~\ref{ch:introduction} and the extent to which they have been answered.

\subsection{Hypothesis Testing}

\begin{table}[h]
\centering
\caption[Hypothesis outcomes summary]{Hypothesis outcomes with key supporting evidence.}
\label{tab:hypothesis_outcomes}
\begin{tabular}{p{0.08\textwidth}p{0.35\textwidth}p{0.15\textwidth}p{0.45\textwidth}}
\toprule
\textbf{H\#} & \textbf{Hypothesis} & \textbf{Outcome} & \textbf{Key Evidence} \\
\midrule
H1 & Dedicated entity channels & \textbf{Partial} & Conv1a specialised, conv2-4a multiplex \\
H2 & SAEs disentangle features & \textbf{Rejected} & SAEs preserve multiplexing structure \\
H3 & Hierarchical preferences & \textbf{Confirmed} & 93\% blue key collection at baseline (cross maze)\\
H4 & SAEs successfully recover performance & \textbf{Confirmed} & 99.6-102\% recovery \\
H5 & Layer specialization & \textbf{Confirmed} & Entity $\rightarrow$ direction transition proven \\
\bottomrule
\end{tabular}
\end{table}

While some of our initial hypotheses proved correct, we were surprised by the discovery of the spatial gating mechanism, where negative activations mark regions of interest and shift systematically as objectives are completed. We also found that SAEs did not work as expected in this case, but as a result were able to validate some of our new results. We were also able to better characterise some of our initial ideas about how the network might work.

\subsection{Research Questions}

\begin{table}[h]
\centering
\caption[Research questions and answers]{Research questions and their answers based on experimental evidence.}
\label{tab:rq_outcomes}
\renewcommand{\arraystretch}{1.5} % Adds 50% more vertical space
\begin{tabular}{p{0.08\textwidth}p{0.35\textwidth}p{0.55\textwidth}}
\toprule
\textbf{RQ} & \textbf{Question} & \textbf{Answer} \\
\midrule
RQ1 & How are goals encoded? & Via spatial gating where negative activations mark regions of interest (Section~\ref{sec:parallel_rollouts},~\ref{sec:negative_activations}) \\
RQ2 & How is the current target determined? & Two-phase: detection at conv1a $\rightarrow$  integration at conv2a-conv4a (Table~\ref{tab:probe_whole_layer}) \\
RQ3 & Can SAEs clarify feature roles? & Yes, but only providing validation (Section~\ref{subsec:sae_structure}) \\
RQ4 & Dedicated vs general channels? & Conv1a specialised conv2a-conv4a generalised (Appendix~\ref{app:channel_probe_results}) \\
RQ5 & Why different steering success rates across entities? & Gem spans in 31/32 channels vs 14 for keys may explain differential steering (Table~\ref{tab:ablation_results}) \\
RQ6 & Do layers have specialised roles? & Yes: entity detection in conv layer, navigation in fc2 (Tables~\ref{tab:probe_whole_layer},~\ref{tab:probe_directional}) \\
\bottomrule
\end{tabular}
\end{table}

\section{Key Findings and Contributions}
\label{sec:contributions}

We primarily demonstrate the unusual and previously undocumented phenomenon of the network using spatial gating to track game state, where negative activations mark regions of interest and shift systematically as objectives are completed. We show this through the clear parallel trajectories of activations as demonstrated by the parallel rollout experiments that have an extremely high correlation (r=0.931 in conv4a), demonstrating that the same channels are active across different game states but at systematically different levels. Most surprisingly, we are able to achieve clear and consistent shifts in behaviour that maintain core network functions with only a simple constant offset to all channel weights, likely by mimicking the activation patterns characteristic of different game stages. This highlights a significant gap in current approaches to MI research, which focuses on identifying which parts of the network are responsible for a certain activation, but fails to consider how temporal patterns in activation levels may encode task progress. 

We further identified a bimodal structure in the network regarding target entity identification early in the network, and a surge in correctly identifying the correct direction to travel later in the network. Entity information crystallizes to 99.3\% probe accuracy in conv3a before collapsing to 20.4\% in fc3, while directional information is not as strong in early layers but strengthens to 67.3\% in fc2. This reveals that convolutional layers determine where things are in the maze before the FC layers determine how to arrive there. This explains why our interventions in the convolutional layers are able to succeed before targets have crystallized.

We also found that SAEs, while of some use in validating our results, failed to provide a great deal of additional insight into the network that would not have been gained otherwise. This highlights the potential pitfalls of trying to apply promising tools before considering simple interventions and experiments that may lead to more insight.

\section{Limitations of the Study}\label{sec:limitations}

This research, while providing several insights, is subject to certain limitations:

Our work is currently limited by the fact that we test only on a single environment, as there are not many environments that have similar qualities such as the extreme generalization encouraged by Procgen, and the staged multi-goal configuration in Heist. To address this, future work could attempt this with MiniGrid environments, though a custom configuration that can establish similar diversity may be needed.

It is also unclear to what extent the intervention channel multiplexing phenomenon generalises to other architectures, and we plan to test this in future work.

Our training methodology differs from standard practice in using unlimited procedurally generated environments rather than the typical 200-500 set of environments one would set. While this was necessary for successful training with our compressed architecture, it may affect the generalization of our findings to models trained under standard Procgen protocols. The unlimited environment diversity likely also contributes to the representational drift we observe, as the network continually trains on new environmental distributions.

We also do not yet fully characterise the function of the intervention spans, or if they translate to an actual unique role within the network's natural functioning. We plan to address this directly through a thorough application of probes that are trained to detect which entity is currently the next entity in order of pursuit by the network. Early work shows that the later convolutional layers do this very successfully, as well as the first fully connected layer. However, to truly understand the role of the intervention spans and navigation channels, it may be necessary to train probes on every channel and every row in the fully connected layer to determine their function.

Intervention experiments, by their nature, can push the network into off-distribution activation states. While this is useful for causal testing, the agent's behaviour under such strong artificial interventions do not represent its typical behaviour. We do not yet have a clear picture of precise circuitry under normal operating conditions. Additionally, the patching interventions were typically only applied with positive activations, but the mean channel activation analysis shown in Figure~\ref{fig:activation_tracking} reveals that for many channels the activations of the network are mostly negative. This means that a large portion of their activity is not being captured accurately in the interventions. We believe future work in understanding the role of the negative activations in the convolutional channels would lead to fruitful insights.

Our approach of simply intervening in every channel with 400 different activation strengths is also an approach that would need to be approached more carefully, due to cost constraints, in larger models.

\section{Future Work}\label{sec:future_work}

This research opens several promising avenues for future investigation. The discovery of spatial gating through activation polarity suggests examining whether similar mechanisms exist in other architectures and domains, particularly in transformer-based models where activation patterns might track different modes of operation.

Testing this mechanism across different RL environments, especially those with hierarchical goal structures or mode-switching requirements, could reveal whether spatial gating is a general principle of efficient neural computation under resource constraints.

The methodological insight that simple mean activation analysis revealed what sophisticated tools missed suggests revisiting other `well-understood' models with basic statistical approaches. This could uncover overlooked organizational principles, particularly in systems previously deemed too polysemantic for interpretation.

The precise control achievable through small activation offsets ($\pm$0.5 at conv2a) warrants investigation into whether training procedures could be designed to encourage such structures, potentially yielding models that are interpretable and controllable by design rather than post-hoc analysis.

Finally, the two-phase architecture's clear separation between `what to target' and `how to navigate' suggests applications in hybrid systems where interpretable goal selection could be combined with powerful but opaque execution modules, potentially offering a path toward more trustworthy autonomous systems.

\section{Final Remarks}
\label{sec:final_remarks}

This thesis demonstrates the practicality of understanding such neural networks, though it also highlights some key bottlenecks in doing this at scale. In particular, many of the tools that have been developed to interpret such models were not able to capture key parts of the network's functioning. This highlights that there are still significant gaps in the way that we approach thinking about neural networks, and in particular networks trained with RL. It also highlights the importance of examining fundamental structures such as activation patterns and high level statistics as a simple first step in approaching such problems.

Given the potential for activation-level encoding to extend beyond this specific environment and architecture, we hope that understanding this particular toy model may alter the way the reader thinks about neural networks, and that cracking the problem of the inner workings of such models may be easier as a result. 
